The Smart Trigger adds a technical authorization step to the AR-15 
without changing how the firearm itself operates. A passive 125kHz RFID
tag worn by the authorized user is read by a custom wound antenna mounted
on the grip. An RDM 6300 module passes the tag ID to the Raspberry Pi Pico W
over UART, and the Pico signals to a micro stepper motor through the SAA1042
driver in order to lock or release the native AR-15 safety selector. Because
the design works with the safety selector that every AR-15 already has, it
can be applied to a wide range of rifles without modifying the lower receiver.

Several of our design decisions came down to balancing ideal conditions with
what we could reliably build on our schedule. A fully integrated custom PCB with
the RP2040 mounted directly would have been the most compact option, but the
risk of hand soldering its micro-leads influenced us to use the Pico W board.
Our PCB implementation then serves as a glorified breadboard that connects 
the Pi Pico W, the RDM6300, the SAA1042, the stepper motor, 9V battery, voltage
regulator, and the antenna all together in a 3-board stack that is compact
enough to fit inside the grip. Designing and printing our own grip in PLA gave
us control over the fit and feel of the final product, and allowed us to iterate
quickly over different designs. For our software design, we built our code on
FreeRTOS to keep RFID processing and motor control as parallel tasks. When a
secure and valid tag is read over UART, we trigger the motor driver
appropriately to unlock the AR-15's safety selector. This is our complete
system design that acheives a supplementory, technical authorization step to
safely operating the AR-15 rifle.